\section{Methodology}

The main econometric exercise estimates the conditional firm-size wage premium
in a log-wage model. The baseline specification is:

\[
  \ln(w_{ist}) =
  \alpha
  + \sum_{g \neq g_0} \beta_g \mathbf{1}\{Size_{ist}=g\}
  + X_{ist}'\gamma
  + \lambda_{s \times t}
  + \varepsilon_{ist},
\]

where \(w_{ist}\) is real hourly labor income for worker \(i\) in sector \(s\) in year \(t\), \(Size_{ist}\) is the harmonized firm-size category, and \(g_0\) is the reference group. We use solo workers as the reference group. This choice makes the coefficient path directly interpretable as the wage premium associated with moving from individual work to increasingly larger firms. Because solo workers may combine self-employment and very small-scale wage employment, we interpret the estimates as conditional gaps relative to the bottom of the firm-size distribution rather than as pure firm effects.

The vector \(X_{ist}\) includes a female-worker indicator, age, age squared,
education dummies, and a formal-employment indicator. The fixed effect
\(\lambda_{s \times t}\)
absorbs sector-year conditions, allowing each broad sector to have its own wage
level in each year. This is important because large firms are not randomly
distributed across sectors, and sectoral shocks may be correlated with both wages
and firm size. The model is estimated with GEIH expansion weights, and standard
errors are clustered by sector.

City fixed effects would be a useful additional robustness check because they
would absorb persistent differences across local labor markets. The harmonized
file used for the regressions does not yet preserve a consistent city or
municipality identifier, but it does retain department/Bogot\'a. Appendix
Table~\ref{tab:firm-size-local-geography-robustness} therefore reports a
stringent local robustness exercise that absorbs sector-department-year fixed
effects.

In the results table, we report four nested specifications. The first column
includes only firm-size indicators. The second adds sector-year fixed effects.
The third adds gender, age, age squared, and education controls, but excludes
formality. The fourth adds formality. Comparing the third and fourth columns
shows how much of the
firm-size gradient is absorbed by differences in formal employment across firm
sizes.

The coefficients \(\beta_g\) are log-point differences relative to solo workers.
For interpretation, we transform them into percentage premiums:

\[
  Premium_g = 100 \times \left(\exp(\hat{\beta}_g)-1\right).
\]

\subsection{Firm-Size Nonlinearities by Formality Status}

The descriptive evidence suggests that the relationship between firm size and
hourly income may differ between formal and informal workers. To allow for this
nonlinearity, we also estimate a specification that interacts all firm-size
categories with formal employment:

\[
  \ln(w_{ist}) =
  \alpha
  + \rho Formal_{ist}
  + \sum_{g \neq g_0} \beta_g \mathbf{1}\{Size_{ist}=g\}
  + \sum_{g \neq g_0} \theta_g
    \mathbf{1}\{Size_{ist}=g\} Formal_{ist}
  + Z_{ist}'\gamma
  + \lambda_{s \times t}
  + \varepsilon_{ist}.
\]

Here \(Z_{ist}\) includes the worker controls other than formality. The
coefficients \(\beta_g\) measure the firm-size premium among informal workers,
relative to informal solo workers. The corresponding premium among formal
workers is \(\beta_g+\theta_g\), relative to formal solo workers. This
specification is deliberately nonparametric in firm size: it does not impose a
linear, quadratic, or monotonic relationship between firm size and wages.

\subsection{Dynamic Premiums and Trend Test}

The same framework can be extended to study the evolution of the premium over
time by interacting firm-size categories with year indicators:

\[
  \ln(w_{ist}) =
  \alpha
  + \sum_{\tau}\sum_{g \neq g_0}
    \beta_{g\tau}\mathbf{1}\{Size_{ist}=g\}\mathbf{1}\{Year_{t}=\tau\}
  + X_{ist}'\gamma
  + \lambda_{s \times t}
  + \varepsilon_{ist}.
\]

This dynamic specification produces a sequence of firm-size premiums by year and
can be used to assess whether the premium has widened, narrowed, or remained
stable. We also estimate a more parsimonious trend specification:

\[
  \ln(w_{ist}) =
  \alpha
  + \sum_{g \neq g_0} \beta_g\mathbf{1}\{Size_{ist}=g\}
  + \sum_{g \neq g_0} \delta_g
    \mathbf{1}\{Size_{ist}=g\}Trend_t
  + X_{ist}'\gamma
  + \lambda_{s \times t}
  + \varepsilon_{ist},
\]

where \(Trend_t\) is a linear year trend normalized to zero in 2008. The
parameter \(\delta_g\) tests whether the wage premium for firm-size group \(g\)
increased over time relative to solo workers. We report one-sided tests of
\(H_0:\delta_g=0\) against \(H_1:\delta_g>0\), because the empirical question is
whether the premium widened over the period.
